\documentclass{article}
\usepackage{ctex}%中文
\begin{document}
\section{A*算法}
算法中有一个启发式函数$h(n)$，它估计从节点$n$到目标节点的代价。A*算法通过结合实际代价$g(n)$和启发式代价$h(n)$来选择下一个要扩展的节点。A*算法的核心思想是优先扩展具有最低$f(n) = g(n) + h(n)$值的节点，其中$g(n)$是从起点到节点$n$的实际代价，$h(n)$是从节点$n$到目标节点的启发式估计代价。A*算法在路径搜索中具有较高的效率和准确性，广泛应用于各种领域，如机器人导航、游戏开发和网络路由等。
Dijkstra算法是一种经典的图算法，用于解决单源最短路径问题。它通过逐步扩展节点来找到从起点到所有其他节点的最短路径。Dijkstra算法的核心思想是使用一个优先队列来存储待扩展的节点，并根据当前已知的最短路径代价进行排序。每次从优先队列中选择具有最低代价的节点进行扩展，并更新其邻居节点的代价。Dijkstra算法保证了在没有负权边的情况下能够找到最短路径，但在存在负权边时可能会出现问题。
贪心最佳优先算法是一种启发式搜索算法，它在每一步选择当前看起来最优的选项，以期望最终能够找到全局最优解。与A*算法不同，贪心最佳优先算法只考虑启发式代价$h(n)$，而不考虑实际代价$g(n)$。因此，贪心最佳优先算法可能会陷入局部最优解，而无法保证找到全局最优解。尽管如此，贪心最佳优先算法在某些情况下仍然可以提供较快的搜索速度，特别是在启发式函数非常准确的情况下。

其中，A*算法在路径搜索中具有较高的效率和准确性，广泛应用于各种领域，如机器人导航、游戏开发和网络路由等。Dijkstra算法是一种经典的图算法，用于解决单源最短路径问题。它通过逐步扩展节点来找到从起点到所有其他节点的最短路径。Dijkstra算法的核心思想是使用一个优先队列来存储待扩展的节点，并根据当前已知的最短路径代价进行排序。每次从优先队列中选择具有最低代价的节点进行扩展，并更新其邻居节点的代价。Dijkstra算法保证了在没有负权边的情况下能够找到最短路径，但在存在负权边时可能会出现问题。贪心最佳优先算法是一种启发式搜索算法，它在每一步选择当前看起来最优的选项，以期望最终能够找到全局最优解。与A*算法不同，贪心最佳优先算法只考虑启发式代价$h(n)$，而不考虑实际代价$g(n)$。因此，贪心最佳优先算法可能会陷入局部最优解，而无法保证找到全局最优解。尽管如此，贪心最佳优先算法在某些情况下仍然可以提供较快的搜索速度，特别是在启发式函数非常准确的情况下。















\end{document}